%%%%%%%%%%%%%%%%%%%%%%%%%%%%%%%%%%%%%%%%%
%
% (c) 2018 by Jennifer Laaser
%
% This work is licensed under the Creative Commons Attribution-NonCommercial-ShareAlike 4.0 International License. To view a copy of this license, visit http://creativecommons.org/licenses/by-nc-sa/4.0/ or send a letter to Creative Commons, PO Box 1866, Mountain View, CA 94042, USA.
%
% The current source for these materials is accessible on Github: https://github.com/jlaaser/pogil-polymers
%
%%%%%%%%%%%%%%%%%%%%%%%%%%%%%%%%%%%%%%%%%

\renewcommand{\figpath}{content/polymchem/stepgrowth/stepgrowth-chemistries/figs}

\begin{activity}[Condensation Polymerizations]

\begin{instructornotes}

	This activity introduces students to key chemistries used for step-growth polymerizations.
	
	After completing this activity, students will be able to:
			\begin{enumerate}
				\item Identify major classes of polymers produced by step-growth polymerization
				\item Determine the polymer produced by a given monomer or monomer pair, and determine the monomer(s) necessary to produce a target polymer
				\item Determine whether or not a reaction qualifies as a condensation polymerization, and if so, identify the small molecule released
			\end{enumerate}
	
			
	\subsection*{Activity summary:}
	\begin{itemize}
		\item \textbf{Activity type:} Learning Cycle
		\item \textbf{Content goals:} Chemistries of step-growth polymerizations
		\item \textbf{Process goals:} %https://pogil.org/uploads/attachments/cj54b5yts006cklx4hh758htf-process-skills-official-pogil-list-2015-original.pdf
			written communication, critical thinking, information processing
		\item \textbf{Duration:} TBD %approx. 45 minutes without class discussion
		\item \textbf{Instructor preparation required:} none beyond knowledge of relevant content
		\item \textbf{Related textbook chapters:}
			\begin{itemize}
				\item \emph{Polymer Chemistry} (Hiemenz \& Lodge): Table 1.2 and sections 2.2.1, 2.5, and 2.6
			\end{itemize}
	\end{itemize}

\end{instructornotes}

	%\textbf{Focus question:} Put a central question for the students to consider through this exercise here.

\begin{model}[Synthesis of a Polyester]

	Esterification reactions are a common type of reaction used to produce polymers by step-growth polymerization.
	In a typical esterification reaction, a hydroxyl group and a carboxylic acid react to form an ester bond:
	
	\centerline{\includegraphics[width=0.7\textwidth]{\figpath/model1_ester-general.pdf}}
	
	One example of a polymerization reaction using this chemistry is the synthesis of poly(6-hydroxycaproic acid) from 6-hydroxycaproic acid monomers:
	
	\centerline{\includegraphics[width=0.9\textwidth]{\figpath/model1_P6HCA.pdf}}

\end{model}


\begin{ctqs}

	\question The 6-hydroxycaproic acid monomer is shown below: \label{ctq:label-6hcpa}
	
	\centerline{\includegraphics[width=0.3\textwidth]{\figpath/model1_6HCA.pdf}}
	
		\begin{enumerate}
			\item As drawn, what type of functional group is on the \emph{left} side of the monomer?
			
				\begin{solution}[0.5in]
				\end{solution}
			
			\item As drawn, what type of functional group is on the \emph{right} side of the monomer?
			
				\begin{solution}[0.5in]
				\end{solution}
		\end{enumerate}
		
\end{ctqs}

\begin{infobox}

	When a monomer used in a step-growth polymerization has different reactive functional groups on each end, it is called an ``AB-type'' monomer.
	
	When a monomer used in a step-growth polymerization has the same reactive functional group on each end, it is called an ``AA-type'' or ``BB-type'' monomer.

\end{infobox}

\begin{ctqs}
		
		\question Would you classify the poly(6-hydroxycaproic acid) monomer used in this synthesis as an AA-type monomer or an AB-type monomer?  Briefly explain your answer in 1-2 complete sentences.
			
				\begin{solution}[1.5in]
				\end{solution}
		
		\question The synthesis of a short 6-hydroxycaproic acid oligomer from three monomers is shown explicitly, below: \label{ctq:6hcpa-oligomer}
	
\vspace{0.25in}	\centerline{\includegraphics[width=0.9\textwidth]{\figpath/model1_3oligomer-explicit.pdf}}
The molecules are color-coded so that you can see which atoms in the oligomer came from which monomer.
		
		\begin{enumerate}
		
			\item Explain, in one or two complete sentences, why you think we classify this polymer as a ``polyester'':
			
				\begin{solution}[2in]
				\end{solution}
		
			\item Explain, in one or two complete sentences, why we generally abbreviate the product of this reaction as
	
	\centerline{\includegraphics[width=0.25\textwidth]{\figpath/model1_3oligomer-shorthand.pdf}}
	
	rather than as
	
	\centerline{\includegraphics[width=0.7\textwidth]{\figpath/model1_incorrect-shorthand.pdf}}
			
				\begin{solution}[1.5in]
				\end{solution}
	
			\item When we write the structure of this oligomer as
	
	\centerline{\includegraphics[width=0.25\textwidth]{\figpath/model1_3oligomer-shorthand.pdf}}
	
	how many monomers make up each repeat unit in the polymer chain?
			
				\begin{solution}[0.5in]
				\end{solution}
			
		\end{enumerate}
	
	\question Consider the following polymer:
	
	\centerline{\includegraphics[width=0.1\textwidth]{\figpath/model1_PMA.pdf}}
	
		\begin{enumerate}
		
			\item Would you be able to produce this polymer by esterification reactions of small molecules?  Why or why not?
			
				\begin{solution}[1.75in]
				\end{solution}
			
			\item Based on your answer to the previous question, would you classify this polymer as a polyester?  Why or why not?
			
				\begin{solution}[1.5in]
				\end{solution}
			
		\end{enumerate}
		
\end{ctqs}
	

\begin{model}[Synthesis of a Polyamide]

Amidation reactions are another type of reaction used to produce polymers by step-growth polymerizations.
For example, acid chlorides can be reacted with primary amines to form an amide bond:
	
	\centerline{\includegraphics[width=0.7\textwidth]{\figpath/model2_amide-general.pdf}}

Commercially, this reaction is used to produce Nomex, a heat-resistant polymer used in oven mitts and firefighters' protective clothing, among other applications.
A reaction scheme for the synthesis of Nomex is shown below:
	
	\centerline{\includegraphics[width=0.9\textwidth]{\figpath/model2_Nomex.pdf}}

\end{model}

\begin{ctqs}
		\question Why is this polymer classified as a polyamide?
			
				\begin{solution}[2in]
				\end{solution}
		
		\question What functional groups does monomer 1 have?   Would you classify this monomer as an AA-type monomer or an AB-type monomer?
			
				\begin{solution}[0.75in]
				\end{solution}
		
		\question What functional groups does monomer 2 have?   Would you classify this monomer as an AA-type monomer or an AB-type monomer?
			
				\begin{solution}[0.75in]
				\end{solution}
		
		\question Explain, in one or two complete sentences, why we might describe this reaction as an ``AA+BB''-type polymerization:
			
				\begin{solution}[1.25in]
				\end{solution}
		
		\question How many monomers make up each repeat unit?
			
				\begin{solution}[0.5in]
				\end{solution}
		
		\question A very similar reaction can be used to make Kevlar, the high-strength polymer used in bulletproof vests and cut-resistant gloves.  Given that Kevlar is produced from the following two monomers,
		
	
	\centerline{\includegraphics[width=0.5\textwidth]{\figpath/model2_Kevlar-monomers.pdf}}
		
		what do you predict the structure of the polymer should look like?
			
				\begin{solution}[1.5in]
				\end{solution}
		
		\question A similar chemistry can also be used to prepare nylon-6,6, a polymer used in many consumer goods.
		The structure of nylon-6,6 is shown below:
		
			\centerline{\includegraphics[width=0.45\textwidth]{\figpath/model2_nylon66.pdf}}
			
		What two monomers would you need to combine to make this polymer?
			
				\begin{solution}[2in]
				\end{solution}
			
\end{ctqs}
	
\begin{infobox}

A polymerization reaction is called a \emph{condensation} polymerization if the reaction produces a small-molecule byproduct that is not part of the polymer chain.

\end{infobox}
	
\begin{ctqs}
		\question Is the esterification reaction in Model 1 a condensation polymerization?  If so, what is the small-molecule byproduct that is produced?
			
				\begin{solution}[1in]
				\end{solution}
		
		\question Is the amidation reaction in Model 2 a condensation polymerization?  If so, what is the small-molecule byproduct that is produced?
			
				\begin{solution}[1in]
				\end{solution}
		
\end{ctqs}

\begin{model}[Other Chemistries used for Step-Growth Polymerization]

Shown below are reactions used to produce a variety of other commercially-important polymers by step-growth polymerization:

\vspace{0.1in}

\textbf{a) polycarbonate} (high transparency and impact resistance; used in DVDs, glasses, etc.)
		
			\centerline{\includegraphics[width=\textwidth]{\figpath/model3_polycarbonate.pdf}}

\vspace{0.1in}

\textbf{b) polyurethanes} (foams; thermoplastic elastomers, e.g. spandex)
		
			\centerline{\includegraphics[width=\textwidth]{\figpath/model3_polyurethane.pdf}}

\vspace{0.1in}

\textbf{c) epoxies} (adhesives; coatings)
		
			\centerline{\includegraphics[width=\textwidth]{\figpath/model3_epoxy.pdf}}

\end{model}

\vspace{0.25in}
\begin{ctqs}

		\question Classify each of the reactions in the above Model as either an ``AB-type'' or ``AA+BB-type'' polymerization:
		
			\begin{enumerate}
				\item ~ \begin{solution}[0.6in]\end{solution}
				\item ~ \begin{solution}[0.6in]\end{solution}
				\item ~ \begin{solution}[0.6in]\end{solution}
			\end{enumerate}
			
		% may need another question in here
		
		\clearpage
		\question Complete the following table for the polymerizations depicted in Models 1-3:
		
			\begin{center}
				\begin{tabular}{|c|c|c|c|c|}
				\hline
					Polymer & ``A'' reactive group & ``B'' reactive group & ``ab'' bond formed & \begin{tabular}{c}Small\\ Molecule\\Byproduct\end{tabular} \\\hline
					Polyester &						
						\includegraphics[width=0.12\textwidth]{\figpath/model3_hydroxyl.pdf} &
						\includegraphics[width=0.15\textwidth]{\figpath/model3_acid.pdf} &
						\includegraphics[width=0.15\textwidth]{\figpath/model3_ester.pdf}&\\\hline
					Polyamide &
						\includegraphics[width=0.12\textwidth]{\figpath/model3_amine.pdf} &
						\includegraphics[width=0.15\textwidth]{\figpath/model3_acidchloride.pdf} &&
						\includegraphics[width=0.08\textwidth]{\figpath/model3_HCl.pdf} \\\hline
					Polycarbonate &
						\includegraphics[width=0.15\textwidth]{\figpath/model3_blank.pdf}&&& \\\hline
					Polyurethane &
						\includegraphics[width=0.15\textwidth]{\figpath/model3_blank.pdf}&&& \\\hline
					Epoxy &
						\includegraphics[width=0.15\textwidth]{\figpath/model3_blank.pdf}&&& \\\hline
				\end{tabular}
			\end{center}
			
		\vspace{0.25in}
		\question Which of the above polymerization reactions would you classify as condensation polymerizations?
			
				\begin{solution}[2in]
				\end{solution}
			
\end{ctqs}
	

\clearpage
\begin{exercises}

		\exercise (other esterification chemistries)
		
		\exercise (difference in polymer structures formed by AB vs AA+BB reactions?)
		
		\exercise (H-bonded structure of polyesters vs. polyamides - justify difference in properties?)
		
		\exercise (phosgene alternatives for polycarbonate synthesis?)
		
		\exercise (epoxidation reaction formed a linear polymer with secondary amines.  however, secondary amines can also attack epoxides.  What type of polymer structure would be formed in this case?)
		
\end{exercises}
	
\end{activity}